% GENERATED FILE. Edit canonical lessons, then run npm run generate:books.
% canonical-source-hash: fnv1a64:fe473542a421cc0b
% canonical-lessons: TA-C385-culnilai, TA-C385-nilaimai, TA-C385-pinnani, TA-C385-cikkal, TA-C385-kaiyal, TA-R385-review-379-385

\chapter{Handling a Situation}
\label{ch:ta-handling-a-situation}
\hlchaptermodality{385}

\begin{chapteropening}
\textbf{By the end of this chapter:} \emph{I can describe a situation and how someone handled it.}

Complete the last lesson of chapter 385: I can describe a situation and how someone handled it.
\end{chapteropening}

\section[\texorpdfstring{\ta{சூழ்நிலை} (cūḻnilai)}{சூழ்நிலை (cūḻnilai)}]{\ta{சூழ்நிலை} (cūḻnilai) — circumstances; the setting}

\label{lesson:TA-C385-culnilai}



\begin{quote}
Before the new one: say the Tamil for at least, then the Tamil for to command.
\end{quote}

\subsection*{You'll want to know: \ta{சூழ்நிலை}}
\textbf{\ta{சூழ்நிலை}} — \emph{cūḻnilai} — \textquotedblleft{}circumstances; the setting\textquotedblright{}.

\begin{quote}
\ta{இந்தச்} \ta{சூழ்நிலையில்} \ta{நான்} \ta{என்ன} \ta{செய்வது}?
\end{quote}

\begin{quote}
\emph{intac cūḻnilaiyil nāṉ eṉṉa ceyvatu?}
\end{quote}

\begin{quote}
What am I to do in these circumstances?
\end{quote}

\ta{எந்தச்} \ta{சூழ்நிலையிலும்}, in any circumstances.

\begin{cousinweb}
\textbf{\ta{சூழ்நிலை}} (\emph{cūḻnilai}) is native: \emph{cūḻ}, \textquotedblleft{}to surround, encompass\textquotedblright{} (DEDR 2698), and \emph{nilai}, \textquotedblleft{}a standing, a state\textquotedblright{} (DEDR 3675): the state of what surrounds us.
\end{cousinweb}

\begin{grammarlens}[title={What you've built}]
One more word for saying what happens if.
\end{grammarlens}

\subsection*{Your turn}
\begin{itemize}
\raggedright
  \item \emph{Say it:} \emph{cūḻnilai}
  \item \emph{Say it:} \emph{cūḻnilai}, once more
  \item \emph{Say it:} \emph{intac cūḻnilaiyil nāṉ eṉṉa ceyvatu?} — What am I to do in these circumstances?
  \item \emph{Recall:} say the Tamil for at least, then the Tamil for to command, then say \emph{cūḻnilai} again
\end{itemize}

\begin{tcolorbox}[breakable,colback=teal!4,colframe=teal!35!black,title={Before you move on}]
What does \ta{சூழ்நிலை} mean? (\textquotedblleft{}Circumstances; the setting\textquotedblright{}.) Say it once more.
\end{tcolorbox}
\section[\texorpdfstring{\ta{நிலைமை} (nilaimai)}{நிலைமை (nilaimai)}]{\ta{நிலைமை} (nilaimai) — the situation; the state of things}

\label{lesson:TA-C385-nilaimai}



\begin{quote}
Before the new one: say the Tamil for to command, then the Tamil for circumstances.
\end{quote}

\subsection*{You'll want to know: \ta{நிலைமை}}
\textbf{\ta{நிலைமை}} — \emph{nilaimai} — \textquotedblleft{}the situation; the state of things\textquotedblright{}.

\begin{quote}
\ta{மழை} \ta{நின்றால்} \ta{நிலைமை} \ta{சரியாகும்}.
\end{quote}

\begin{quote}
\emph{maḻai niṉṟāl nilaimai cariyākum.}
\end{quote}

\begin{quote}
If the rain stops, the situation will improve.
\end{quote}

\ta{நிலைமை} \ta{மோசமாக} \ta{இருக்கிறது}, the situation is bad.

\begin{cousinweb}
\textbf{\ta{நிலைமை}} (\emph{nilaimai}) is native Dravidian: DEDR 3675 lists Tamil \emph{nil}, \textquotedblleft{}to stand, stay\textquotedblright{}, and \emph{nilaimai}, \textquotedblleft{}state, condition\textquotedblright{}.
\end{cousinweb}

\begin{grammarlens}[title={What you've built}]
One more word for saying what happens if.
\end{grammarlens}

\subsection*{Your turn}
\begin{itemize}
\raggedright
  \item \emph{Say it:} \emph{nilaimai}
  \item \emph{Say it:} \emph{nilaimai}, once more
  \item \emph{Say it:} \emph{maḻai niṉṟāl nilaimai cariyākum.} — If the rain stops, the situation will improve.
  \item \emph{Recall:} say the Tamil for to command, then the Tamil for circumstances, then say \emph{nilaimai} again
\end{itemize}

\begin{tcolorbox}[breakable,colback=teal!4,colframe=teal!35!black,title={Before you move on}]
What does \ta{நிலைமை} mean? (\textquotedblleft{}The situation; the state of things\textquotedblright{}.) Say it once more.
\end{tcolorbox}
\section[\texorpdfstring{\ta{பின்னணி} (piṉṉaṇi)}{பின்னணி (piṉṉaṇi)}]{\ta{பின்னணி} (piṉṉaṇi) — a background}

\label{lesson:TA-C385-pinnani}



\begin{quote}
Before the new one: say the Tamil for circumstances, then the Tamil for the situation.
\end{quote}

\subsection*{You'll want to know: \ta{பின்னணி}}
\textbf{\ta{பின்னணி}} — \emph{piṉṉaṇi} — \textquotedblleft{}a background\textquotedblright{}.

\begin{quote}
\ta{இந்தக்} \ta{கதையின்} \ta{பின்னணி} \ta{ஒரு} \ta{சிறிய} \ta{கிராமம்}.
\end{quote}

\begin{quote}
\emph{intak kataiyiṉ piṉṉaṇi oru ciṟiya kirāmam.}
\end{quote}

\begin{quote}
The background of this story is a small village.
\end{quote}

Also background music in a film: \ta{பின்னணி} \ta{இசை}.

\begin{cousinweb}
\textbf{\ta{பின்னணி}} (\emph{piṉṉaṇi}) is native: \emph{piṉ}, \textquotedblleft{}behind, after\textquotedblright{} (DEDR 4205), and \emph{aṇi}, \textquotedblleft{}array, row\textquotedblright{} (DEDR 117): the rank at the back.
\end{cousinweb}

\begin{grammarlens}[title={What you've built}]
One more word for saying what happens if.
\end{grammarlens}

\subsection*{Your turn}
\begin{itemize}
\raggedright
  \item \emph{Say it:} \emph{piṉṉaṇi}
  \item \emph{Say it:} \emph{piṉṉaṇi}, once more
  \item \emph{Say it:} \emph{intak kataiyiṉ piṉṉaṇi oru ciṟiya kirāmam.} — The background of this story is a small village.
  \item \emph{Recall:} say the Tamil for circumstances, then the Tamil for the situation, then say \emph{piṉṉaṇi} again
\end{itemize}

\begin{tcolorbox}[breakable,colback=teal!4,colframe=teal!35!black,title={Before you move on}]
What does \ta{பின்னணி} mean? (\textquotedblleft{}A background\textquotedblright{}.) Say it once more.
\end{tcolorbox}
\section[\texorpdfstring{\ta{சிக்கல்} (cikkal)}{சிக்கல் (cikkal)}]{\ta{சிக்கல்} (cikkal) — a complication; a snag}

\label{lesson:TA-C385-cikkal}



\begin{quote}
Before the new one: say the Tamil for the situation, then the Tamil for a background.
\end{quote}

\subsection*{You'll want to know: \ta{சிக்கல்}}
\textbf{\ta{சிக்கல்}} — \emph{cikkal} — \textquotedblleft{}a complication; a snag\textquotedblright{}.

\begin{quote}
\ta{ஆவணங்கள்} \ta{இல்லாவிட்டால்} \ta{சிக்கல்} \ta{வரும்}.
\end{quote}

\begin{quote}
\emph{āvaṇaṅkaḷ illāviṭṭāl cikkal varum.}
\end{quote}

\begin{quote}
If the documents are missing, there will be complications.
\end{quote}

\ta{சிக்கலான}, complicated: \ta{சிக்கலான} \ta{கேள்வி}, a tricky question.

\begin{cousinweb}
\textbf{\ta{சிக்கல்}} (\emph{cikkal}) is native Dravidian: DEDR 2498 lists Tamil \emph{cikku}, \textquotedblleft{}to become entangled, complicated, be caught; a tangle\textquotedblright{}; \emph{cikkal} is its noun.
\end{cousinweb}

\begin{grammarlens}[title={What you've built}]
One more word for saying what happens if.
\end{grammarlens}

\subsection*{Your turn}
\begin{itemize}
\raggedright
  \item \emph{Say it:} \emph{cikkal}
  \item \emph{Say it:} \emph{cikkal}, once more
  \item \emph{Say it:} \emph{āvaṇaṅkaḷ illāviṭṭāl cikkal varum.} — If the documents are missing, there will be complications.
  \item \emph{Recall:} say the Tamil for the situation, then the Tamil for a background, then say \emph{cikkal} again
\end{itemize}

\begin{tcolorbox}[breakable,colback=teal!4,colframe=teal!35!black,title={Before you move on}]
What does \ta{சிக்கல்} mean? (\textquotedblleft{}A complication; a snag\textquotedblright{}.) Say it once more.
\end{tcolorbox}
\section[\texorpdfstring{\ta{கையாள்} (kaiyāḷ)}{கையாள் (kaiyāḷ)}]{\ta{கையாள்} (kaiyāḷ) — to handle; to deal with}

\label{lesson:TA-C385-kaiyal}



\begin{quote}
Before the new one: say the Tamil for a background, then the Tamil for a complication.
\end{quote}

\subsection*{You'll want to know: \ta{கையாள்}}
\textbf{\ta{கையாள்}} — \emph{kaiyāḷ} — \textquotedblleft{}to handle; to deal with\textquotedblright{}.

\begin{quote}
\ta{அவர்} \ta{இந்தச்} \ta{சிக்கலை} \ta{அமைதியாகக்} \ta{கையாண்டார்}.
\end{quote}

\begin{quote}
\emph{avar intac cikkalai amaitiyākak kaiyāṇṭār.}
\end{quote}

\begin{quote}
He handled this complication calmly.
\end{quote}

Past (I …): \ta{கையாண்டேன்}, \emph{kaiyāṇṭēṉ}.

What is handled takes -\ta{ஐ}. The past \ta{கையாண்டேன்} is like \ta{ஆண்டேன்} from \ta{ஆள்}.

\begin{cousinweb}
\textbf{\ta{கையாள்}} (\emph{kaiyāḷ}) is native: \emph{kai}, \textquotedblleft{}hand\textquotedblright{} (DEDR 2023), and \emph{āḷ}, \textquotedblleft{}to rule, manage, use\textquotedblright{} (DEDR 5157): to manage with the hand.
\end{cousinweb}

\begin{grammarlens}[title={What you've built}]
One more word for saying what happens if.
\end{grammarlens}

\subsection*{Your turn}
\begin{itemize}
\raggedright
  \item \emph{Say it:} \emph{kaiyāḷ}
  \item \emph{Say it:} \emph{kaiyāḷ}, once more
  \item \emph{Say it:} \emph{avar intac cikkalai amaitiyākak kaiyāṇṭār.} — He handled this complication calmly.
  \item \emph{Recall:} say the Tamil for a background, then the Tamil for a complication, then say \emph{kaiyāḷ} again
\end{itemize}

\begin{tcolorbox}[breakable,colback=teal!4,colframe=teal!35!black,title={Before you move on}]
What does \ta{கையாள்} mean? (\textquotedblleft{}To handle; to deal with\textquotedblright{}.) Say it once more.
\end{tcolorbox}
\section[(dialogue)]{Looking back: from Helping Others to Handling a Situation}

\label{lesson:TA-R385-review-379-385}



\begin{quote}
Say the Tamil for a volunteer and a beneficiary.
\end{quote}

\subsection*{Your turn}
Say these aloud:

\begin{itemize}
\raggedright
  \item a volunteer, a beneficiary, funds, to muster, welfare
  \item responsibility, an obligation, to dedicate, to contribute, to cooperate
  \item a nation, a refugee, a symbol, to unite, a representative
  \item a condition, to be subject to, eligibility, an exception, if need be
  \item a possibility, probability, probably, an occasion, or else
  \item necessity, compulsion, a limit, at least, to command
  \item circumstances, the situation, a background, a complication, to handle
\end{itemize}

\begin{tcolorbox}[breakable,colback=teal!4,colframe=teal!35!black,title={Before you move on}]
A volunteer? (\textbf{\ta{தன்னார்வலர்}}, \emph{taṉṉārvalar}.) A beneficiary? (\textbf{\ta{பயனாளி}}, \emph{payaṉāḷi}.) Funds? (\textbf{\ta{நிதி}}, \emph{niti}.) To muster? (\textbf{\ta{திரட்டு}}, \emph{tiraṭṭu}.) Welfare? (\textbf{\ta{நலன்}}, \emph{nalaṉ}.) Responsibility? (\textbf{\ta{பொறுப்பு}}, \emph{poṟuppu}.) An obligation? (\textbf{\ta{கடப்பாடு}}, \emph{kaṭappāṭu}.) To dedicate? (\textbf{\ta{அர்ப்பணி}}, \emph{arppaṇi}.) To contribute? (\textbf{\ta{பங்களி}}, \emph{paṅkaḷi}.) To cooperate? (\textbf{\ta{ஒத்துழை}}, \emph{ottuḻai}.) A nation? (\textbf{\ta{தேசம்}}, \emph{tēcam}.) A refugee? (\textbf{\ta{அகதி}}, \emph{akati}.) A symbol? (\textbf{\ta{சின்னம்}}, \emph{ciṉṉam}.) To unite? (\textbf{\ta{ஒன்றிணை}}, \emph{oṉṟiṇai}.) A representative? (\textbf{\ta{பிரதிநிதி}}, \emph{piratiniti}.) A condition? (\textbf{\ta{நிபந்தனை}}, \emph{nipantaṉai}.) To be subject to? (\textbf{\ta{உட்படு}}, \emph{uṭpaṭu}.) Eligibility? (\textbf{\ta{தகுதி}}, \emph{takuti}.) An exception? (\textbf{\ta{விதிவிலக்கு}}, \emph{vitivilakku}.) If need be? (\textbf{\ta{வேண்டுமானால்}}, \emph{vēṇṭumāṉāl}.) A possibility? (\textbf{\ta{சாத்தியம்}}, \emph{cāttiyam}.) Probability? (\textbf{\ta{நிகழ்தகவு}}, \emph{nikaḻtakavu}.) Probably? (\textbf{\ta{அநேகமாக}}, \emph{anēkamāka}.) An occasion? (\textbf{\ta{சந்தர்ப்பம்}}, \emph{cantarppam}.) Or else? (\textbf{\ta{இல்லையேல்}}, \emph{illaiyēl}.) Necessity? (\textbf{\ta{அவசியம்}}, \emph{avaciyam}.) Compulsion? (\textbf{\ta{கட்டாயம்}}, \emph{kaṭṭāyam}.) A limit? (\textbf{\ta{வரம்பு}}, \emph{varampu}.) At least? (\textbf{\ta{குறைந்தபட்சம்}}, \emph{kuṟaintapaṭcam}.) To command? (\textbf{\ta{ஆணையிடு}}, \emph{āṇaiyiṭu}.) Circumstances? (\textbf{\ta{சூழ்நிலை}}, \emph{cūḻnilai}.) The situation? (\textbf{\ta{நிலைமை}}, \emph{nilaimai}.) A background? (\textbf{\ta{பின்னணி}}, \emph{piṉṉaṇi}.) A complication? (\textbf{\ta{சிக்கல்}}, \emph{cikkal}.) To handle? (\textbf{\ta{கையாள்}}, \emph{kaiyāḷ}.)
\end{tcolorbox}
